\chapter{Zahl und Vorzeichen}

\section{Die Zahl zwischen beiden}

Aristoteles hat die Zeit die Zahl der Bewegung nach dem Früher und Später genannt. Fink liest die Definition als Fundierung, Größe, Bewegung, Zeit, und bemerkt, dass die Zahl die Bewegung \zitat{umfasst}, wie der Ort als Gefäß den Körper umfasst.\footnote{Fink, Frühgeschichte, Kapitel zur aristotelischen Zeit-Analytik.} Die Zeit ist also bei Aristoteles nicht selbst eine Größe, sondern das, womit eine Größe, die Bewegung, gezählt wird. Damit steht die Zahl zwischen Zeit und Raum: Sie zählt das Räumliche, die durchlaufene Größe, und sie ist das Zeitliche, das Früher und Später.

Hegel hat die Zahl zwischen das Sinnliche und den Gedanken gestellt, und die Arithmetik, die man der Zeit als ihre Wissenschaft zuordnen möchte, reduziert die Zeit auf das tote Eins. Das Eins ist das Jetzt, aus dem das Werden herausgenommen ist; eine Zeit aus Einsen ist keine Zeit.\footnote{Bonsiepen, Hegels Raum-Zeit-Lehre, zur Frage einer Wissenschaft der Zeit.}

Baumann hat die gewöhnliche Ableitung der Zahl aus der Zeit, die sich auf die aristotelische Definition beruft, dreifach bestritten: \zitat{Erstens ist die Zahl discret, die Zeit continuirlich; sodann ist das blosse Nacheinander noch keine Zahl, so wenig wie es an sich schon Zeit ist, aber was in beiden Fällen hinzukommen muss, ist jedesmal etwas Anderes [\ldots]; endlich ist der Begriff der Einheit denkbar auch ohne die Zeit.}\footnote{Baumann, Lehren, Bd.\,2, S.\,668.} Das Eins ist ein geistiger Begriff, anwendbar wo und wie wir wollen; das Zusammenfassen von Eins und Eins zu Zwei ist ein geistiges Tun, das auch Gleichzeitiges betreffen kann. Die Zahl ist also weder aus der Zeit noch aus dem Raum, sondern aus dem Geist; und wo sie auf Zeit und Raum angewandt wird, verhält sie sich zu beiden verschieden. Den Raum lässt sie bestehen, denn seine Teile sind nebeneinander und lassen sich zählen, ohne dass einer vergeht. Die Zeit tötet sie, denn ihre Teile sind nicht nebeneinander, und wer sie zählt, hat sie nebeneinandergelegt. Das ist der Punkt, an dem die Parallelisierung entsteht: nicht in der Zeit und nicht im Raum, sondern in der Zahl, die beide gleich behandelt, weil sie beiden gleich fremd ist.

\section{Koordinaten und Zuordnung}

Reichenbach hat die Zahl in der Physik an ihren Ort gewiesen. Die Koordinaten sind bloße Zahlen; was sie bedeuten, wird durch Zuordnungsdefinitionen festgelegt, die sagen, welcher Körper starr ist, welche Uhr gleichförmig geht, welches Signal die Gleichzeitigkeit definiert. Die Zahl ist nicht das Maß, sie ist die Stelle, an die das Maß geschrieben wird. Schneider hat aus der allgemeinen Relativitätstheorie dasselbe gelesen: Die Koordinaten sind Symbole, messbar ist das Linienelement, und nur durch die Materie kommt physikalische Bedeutung hinein.\footnote{Reichenbach, Raum-Zeit-Lehre, \S\,4 und \S\,35; Schneider, Raum-Zeit-Problem, Kapitel zur allgemeinen Relativitätstheorie.}

Das heißt: Die Gleichbehandlung von Zeit und Raum, die in der Koordinatenmannigfaltigkeit vorliegt, ist eine Gleichbehandlung von Zahlen. Vier Zahlen legen ein Ereignis fest, das ist alles, was die Dimensionszahl zunächst besagt. Was die Zahlen bedeuten, ob sie Längen sind oder Zeiten, wird nicht von der Mannigfaltigkeit gesagt, sondern von der Zuordnung, und die Zuordnung ist für die eine Koordinate eine andere als für die drei übrigen.

\section{Das Vorzeichen}

Wo die Zuordnung anders ist, muss die Mannigfaltigkeit es irgendwo anzeigen. Sie zeigt es im Vorzeichen. Das Linienelement der Raum-Zeit ist keine Summe von Quadraten, sondern eine Differenz: Das Zeitglied hat das entgegengesetzte Vorzeichen der Raumglieder. Minkowski hat das Vorzeichen versteckt, indem er die Zeitkoordinate mit der imaginären Einheit multiplizierte; dann sieht das Linienelement aus wie eine euklidische Länge, und die Mannigfaltigkeit wie ein Raum mit einer Dimension mehr. Schneider referiert den Grund: Minkowski führte die Zeit als imaginäre Koordinate ein, um Symmetrie in der quadratischen Form zu erreichen. Die Symmetrie ist eine Symmetrie der Form; die Sache lässt sie, wie sie ist.\footnote{Schneider, Raum-Zeit-Problem, Kapitel zur speziellen Relativitätstheorie.}

Reichenbach nimmt das Vorzeichen ernst. Es ist für ihn der mathematische Ausdruck der Sonderstellung der Zeit, und die Sonderstellung ist die zuvor beschriebene: Die Zeit ist die Faserrichtung, die Dimension der Wirkungslinien, und nur sie. Dass man dieses Vorzeichen vor das Zeitglied oder, mit umgekehrter Konvention, vor die drei Raumglieder schreibt, ist Wahl. Dass es eines gibt, das ein Glied von drei anderen trennt, ist keine Wahl.\footnote{Reichenbach, Raum-Zeit-Lehre, \S\,27, \S\,41 und \S\,43.}

Hier ist eine Unterscheidung nötig, die leicht verwischt wird. Die Signatur der Metrik ist nicht die Asymmetrie. Sie ist deren metrische Spur, das, was von der Asymmetrie übrigbleibt, wenn man Zeit und Raum als Koordinaten schreibt. Die Asymmetrie selbst liegt vor der Metrik: in der Ordnung, in der Grenze, in der Richtung, im Modus. Das Vorzeichen zeigt an, dass diese Asymmetrie da ist; es sagt nicht, worin sie besteht. Wer die Signatur liest, als sagte sie, die Zeit sei reell und der Raum imaginär oder umgekehrt, liest eine Konvention als Befund. Was sich lesen lässt, ist nur dies: Eine Koordinate ist anders als drei.

Brentano hat die Rede von der Zeit als weiterer Raumdimension eine \zitat{unschädliche Fiktion} genannt.\footnote{Brentano, Untersuchungen, S.\,209.} Das trifft die Lage genau. Die Fiktion ist unschädlich, solange man weiß, dass sie eine ist; sie ist in der Physik unentbehrlich, weil die Physik die Sache will und der Modus in der Sache nicht vorkommt. Sie wird schädlich in dem Augenblick, in dem man das Vorzeichen für eine Auskunft darüber hält, was Zeit und Raum sind. Es ist die Auskunft, dass sie nicht dasselbe sind, und sonst nichts.

\section{Was die Zahl nicht erreicht}

Was die Zahl an der Zeit nicht erreicht, hat Brentano gesagt: den Modus, das Vergangen- und Künftigsein, das keine Bestimmung des Gegenstandes ist. Was sie am Raum nicht erreicht, hat Guzzoni gesagt: die Weite, die keine Höhe, Breite oder Länge ist. In beiden Fällen bleibt ein Unmaß, das dem Maß vorausgeht. Aber die beiden Unmaße sind nicht gleich. Die Weite ist der Raum, wie er ist, bevor man ihn misst; sie ist Raum. Der Modus ist die Zeit, wie sie ist, bevor man sie zählt; und er ist nicht Zeit im Sinne der Linie, sondern das, was die Linie zur Zeit macht. Die Zahl lässt vom Raum die Weite zurück und von der Zeit die Zeit.
